\chapter{引力 (Gravitational Forces)}

大多数航天器动力学问题的求解都需要考虑引力和力矩。本章以牛顿万有引力定律为出发点，针对这一背景特别开发相关力和力矩的表达式。

\section{两个粒子的引力相互作用 (Gravitational Interaction of Two Particles)}

\begin{mydefinition}{引力定律}
质量为 \(m\) 的粒子 \(P\) 在质量为 \(\overline{m}\) 的粒子 \(\overline{P}\) 存在时会受到一个力 \(\mathbf{F}\) 的作用。该力沿着连接 \(P\) 与 \(\overline{P}\) 的直线，方向从 \(P\) 指向 \(\overline{P}\)，大小与 \(m\) 和 \(\overline{m}\) 的乘积成正比，与 \(P\) 和 \(\overline{P}\) 之间距离的平方成反比。
因此，如果 \(\mathbf{p}\) 是从 \(\overline{P}\) 指向 \(P\) 的位置矢量，则：
\begin{equation}
\mathbf{F} = -G\overline{m}m\mathbf{p}(\mathbf{p}^2)^{-3/2} \tag{1}
\end{equation}
其中 \(G \approx 6.6732 \times 10^{-11} \ \text{N}\cdot\text{m}^2\text{kg}^{-2}\) 为万有引力常数。
\end{mydefinition}

根据牛顿第三定律，粒子 \(P\) 对 \(\overline{P}\) 的引力 \(\overline{\mathbf{F}}\) 满足：
\begin{equation}
\overline{\mathbf{F}} = -\mathbf{F} \tag{2}
\end{equation}

\section{粒子对物体施加的力 (Force Exerted on a Body by a Particle)}

\begin{mydefinition}{等效引力}
由质量为 \(\overline{m}\) 的粒子 \(\overline{P}\) 对物体 \(B\) 的各质点施加的引力系等效于一个作用于 \(\overline{P}\) 且作用线通过 \(\overline{P}\) 的合力 \(\mathbf{F}\)。注意，该作用线不一定通过 \(B\) 的质心 \(B^*\)。
\end{mydefinition}

如果 \(B\) 由 \(N\) 个质量为 \(m_i\) 的粒子 \(P_i\) 组成，且 \(\mathbf{p}_i\) 是 \(P_i\) 相对于 \(\overline{P}\) 的位置矢量，则：
\begin{equation}
\mathbf{F} = -G\overline{m}\sum_{i=1}^N m_i\mathbf{p}_i(\mathbf{p}_i^2)^{-3/2} \tag{1}
\end{equation}
如果 \(B\) 是连续的介质，\(\rho\) 是 \(B\) 中任意点 \(P\) 的质量密度，\(\mathbf{p}\) 是 \(P\) 相对于 \(\overline{P}\) 的位矢，\(d\tau\) 是元体积（或面积/长度），则：
\begin{equation}
\mathbf{F} = -G\overline{m}\int \mathbf{p}(\mathbf{p}^2)^{-3/2}\rho d\tau \tag{2}
\end{equation}

\begin{mydefinition}{引力中心}
在合力 \(\mathbf{F}\) 的作用线上总可以找到一个点 \(B'\)，使得放置在 \(B'\) 处、质量等于 \(B\) 的质点所受到的来自 \(\overline{P}\) 的引力等于 \(\mathbf{F}\)。该点 \(B'\) 称为 \(B\) 对吸引粒子 \(\overline{P}\) 的**引力中心**。引力中心不一定与质心重合。
\end{mydefinition}

\begin{proof}[等效力系统的证明]
\(\overline{P}\) 对 \(B\) 每一微元施加的力 \(d\mathbf{F}\) 沿连线 \(P\to \overline{P}\) 作用。系统合力为 \(\int d\mathbf{F}\)，关于点 \(\overline{P}\) 的力矩为零。因此，根据力系等价定义，这个合力 \(\mathbf{F}\) 等价于原引力系统。
\end{proof}

\section{粒子对小物体施加的力 (Force Exerted on a Small Body by a Particle)}

\begin{myTheorem}{小物体受力展开}
当物体 \(B\) 的质心 \(B^*\) 与粒子 \(\overline{P}\) 之间的距离 \(R\) 大于 \(B\) 上任意点到 \(B^*\) 的最大距离时，式 (2.2.2) 中的力 \(\mathbf{F}\) 可展开为：
\begin{equation}
\mathbf{F} = -\frac{G\overline{m}m}{R^2}\left[\mathbf{a}_1 + \sum_{i=2}^{\infty}\mathbf{f}^{(i)}\right] \tag{1}
\end{equation}
其中 \(m\) 是 \(B\) 的质量，\(\mathbf{a}_1\) 是从 \(\overline{P}\) 指向 \(B^*\) 的单位向量，使得 \(\mathbf{R} = R\mathbf{a}_1\)。\(\mathbf{f}^{(i)}\) 是 \(|\mathbf{r}|/R\) 的 \(i\) 阶项组合。
\end{myTheorem}
特别地：
\begin{equation}
\mathbf{f}^{(2)} = \frac{1}{mR^2}\left\{\frac{3}{2}\left[\mathrm{tr}(\mathbf{I}) - 5\mathbf{a}_1\cdot \mathbf{I}\cdot \mathbf{a}_1\right]\mathbf{a}_1 + 3\mathbf{I}\cdot \mathbf{a}_1\right\} \tag{3}
\end{equation}
其中 \(\mathbf{I}\) 是 \(B\) 的**中心惯性并矢**，\(\mathrm{tr}(\mathbf{I})\) 是它的迹。

这给出了两个实用近似：
\begin{equation}
\mathbf{F}\approx \hat{\mathbf{F}}\triangleq -\frac{G\overline{m}m}{R^2}\mathbf{a}_1 \tag{5}
\end{equation}
以及：
\begin{equation}
\mathbf{F}\approx \widetilde{\mathbf{F}}\triangleq -\frac{G\overline{m}m}{R^2}[\mathbf{a}_1 + \mathbf{f}^{(2)}] \tag{6}
\end{equation}

\begin{proof}[级数展开的推导]
将 \(\mathbf{p} = \mathbf{R} + \mathbf{r}\) 代入 (2.2.2)，令 \(\mathbf{q} = \mathbf{r}/R\)，对分母 \((1 + 2\mathbf{a}_1\cdot\mathbf{q} + \mathbf{q}^2)^{-3/2}\) 应用二项式级数展开，并利用 \(B^*\) 是质心的性质 \(\int \mathbf{q}\rho d\tau = 0\) 即可得到上述结果。
\end{proof}

\section{小物体对小物体施加的力 (Force Exerted on a Small Body by a Small Body)}

设 \(B\) 和 \(\overline{B}\) 为两个物体，其质心 \(B^*\) 和 \(\overline{B}^*\) 间的距离 \(R\) 远大于各自的尺寸。则 \(\overline{B}\) 对 \(B\) 的引力合力为：
\begin{equation}
\mathbf{F} = -\frac{G\overline{m}m}{R^2}\left[\mathbf{a}_1 + \sum_{i=2}^{\infty}\mathbf{f}^{(i)} + \sum_{i=2}^{\infty}\overline{\mathbf{f}}^{(i)} + \sum_{i=2}^{\infty}\sum_{j=2}^{\infty}\mathbf{f}^{(ij)}\right] \tag{1}
\end{equation}
\(\mathbf{f}^{(i)}\) 和 \(\overline{\mathbf{f}}^{(i)}\) 是仅与 \(B\) 和仅与 \(\overline{B}\) 相关的惯性项，其二阶项为：
\begin{align}
\mathbf{f}^{(2)} &= \frac{1}{mR^2}\left\{\frac{3}{2}\left[\mathrm{tr}(\mathbf{I}) - 5\mathbf{a}_1\cdot \mathbf{I}\cdot \mathbf{a}_1\right]\mathbf{a}_1 + 3\mathbf{I}\cdot \mathbf{a}_1\right\} \tag{2} \\
\overline{\mathbf{f}}^{(2)} &= \frac{1}{\overline{m}R^2}\left\{\frac{3}{2}\left[\mathrm{tr}(\overline{\mathbf{I}}) - 5\mathbf{a}_1\cdot \overline{\mathbf{I}}\cdot \mathbf{a}_1\right]\mathbf{a}_1 + 3\overline{\mathbf{I}}\cdot \mathbf{a}_1\right\} \tag{3}
\end{align}

\begin{proof}[推导]
利用 2.3 节的结果代表 \(\overline{B}\) 上微元对 \(B\) 的引力，并对 \(\overline{B}\) 积分。通过对被积函数的级数展开和重组，即可得到 \(\mathbf{f}^{(2)}\) 和 \(\overline{\mathbf{f}}^{(2)}\)。
\end{proof}

\section{质心体 (Centrobaric Bodies)}

\begin{mydefinition}{质心体}
满足如下条件的物体 \(B\) 称为**质心体**：对于任何粒子 \(\overline{P}\)，作用在 \(B\) 上的引力 \(\mathbf{F}\) 总是等于假设将其所有质量集中于质心 \(B^*\) 时的引力：
\begin{equation}
\mathbf{F} = -\frac{G\overline{m}m}{R^2}\mathbf{a}_1 \tag{1}
\end{equation}
换言之，这种物体的引力中心和质心在任何位置都重合。
\end{mydefinition}

\begin{myTheorem}{质心体的惯性性质}
质心体必具有以下性质：关于通过其质心 \(B^*\) 的任何直线的转动惯量都相等。换言之，其中心惯性椭球为一个球。
\end{myTheorem}
\begin{proof}
对于质心体，式 (1) 成立意味着 \(\sum_{i=2}^\infty \mathbf{f}^{(i)} = 0\)。由各项阶次互不相关，必有 \(\mathbf{f}^{(2)} = 0\)。将 \(\mathbf{f}^{(2)}\) 的表达式代入，可得 \(\mathbf{a}_1 \cdot \mathbf{I} \cdot \mathbf{a}_1 = \frac{1}{3}\mathrm{tr}(\mathbf{I})\)。由于 \(\mathrm{tr}(\mathbf{I})\) 与方向无关，故任何方向的转动惯量均相等。
\end{proof}

\section{粒子对物体施加的力矩 (Moment Exerted on a Body by a Particle)}

\begin{myTheorem}{引力矩的近似}
由粒子 \(\overline{P}\) 作用在物体 \(B\) 上产生的关于其质心 \(B^*\) 的力矩 \(\mathbf{M}\) 等于：
\begin{equation}
\mathbf{M} = -\mathbf{R}\times\mathbf{F} \tag{1}
\end{equation}
其中 \(\mathbf{F}\) 是 2.2 节中给出的合力，\(\mathbf{R}\) 是 \(\overline{P}\) 指向 \(B^*\) 的位矢。
当 \(R\) 远大于 \(B\) 的尺寸时，力矩可近似为：
\begin{equation}
\mathbf{M} \approx \widetilde{\mathbf{M}} \triangleq \frac{3G\overline{m}}{R^3}\mathbf{a}_1\times\mathbf{I}\cdot\mathbf{a}_1 \tag{3}
\end{equation}
\end{myTheorem}

利用惯性矩分量 \(I_{jk}\) 或主转动惯量 \(I_1,I_2,I_3\)，可写出：
\begin{equation}
\widetilde{\mathbf{M}} = \frac{3G\overline{m}}{R^3}(I_{21}\mathbf{a}_3 - I_{31}\mathbf{a}_2) \tag{5}
\end{equation}
或：
\begin{equation}
\widetilde{\mathbf{M}} = \frac{3G\overline{m}}{R^3} \left[ \mathbf{b}_1 (I_3 - I_2)C_{12}C_{13} + \mathbf{b}_2 (I_1 - I_3)C_{13}C_{11} + \mathbf{b}_3 (I_2 - I_1)C_{11}C_{12} \right] \tag{8}
\end{equation}
\begin{proof}
利用式 (1) 和 2.3 节中 \(\mathbf{F}\) 的展开式，保留最低阶项即可得证。
\end{proof}

\section{小物体对小物体施加的力矩 (Moment Exerted on a Small Body by a Small Body)}

当 \(B\) 和 \(\overline{B}\) 之间的距离远大于它们各自尺寸时，\(\overline{B}\) 对 \(B\) 的引力系统关于 \(B^*\) 的力矩为：
\begin{equation}
\mathbf{M} \approx \widetilde{\mathbf{M}} \triangleq \frac{3G\overline{m}}{R^3}\mathbf{a}_1\times\mathbf{I}\cdot\mathbf{a}_1 \tag{2}
\end{equation}

\section{邻近物体 (Proximate Bodies)}

在处理非常靠近的两个物体的引力相互作用时，尽管通常假设 \(R\) 很大以适用展开式，但式 (2.7.2) 往往在近距离依然非常有效，只要 \(\overline{B}\) 接近质心体，且 \(R\) 相对 \(B\) 的尺寸依然足够大。此时，\(\overline{B}\) 几乎像点质量一样作用于 \(B\)。

\section{对向量求微分 (Differentiation with Respect to a Vector)}

\begin{mydefinition}{向量梯度、散度与拉普拉斯}
若标量 \(F\) 依赖于向量 \(\mathbf{v}\)，则定义：
\begin{equation}
\nabla_v F \triangleq \frac{\partial F}{\partial v_1}\mathbf{a}_1 + \frac{\partial F}{\partial v_2}\mathbf{a}_2 + \frac{\partial F}{\partial v_3}\mathbf{a}_3 \tag{1}
\end{equation}
其中 \(v_i = \mathbf{v}\cdot\mathbf{a}_i\)。这称为标量 \(F\) 对向量 \(\mathbf{v}\) 的梯度。
类似的，对于向量 \(\mathbf{F}\)，定义散度和梯度张量：
\begin{align}
\nabla_v \cdot \mathbf{F} &\triangleq \frac{\partial F_1}{\partial v_1} + \frac{\partial F_2}{\partial v_2} + \frac{\partial F_3}{\partial v_3} \tag{2} \\
\nabla_v \mathbf{F} &\triangleq (\nabla_v F_1)\mathbf{a}_1 + (\nabla_v F_2)\mathbf{a}_2 + (\nabla_v F_3)\mathbf{a}_3 \tag{3}
\end{align}
当 \(\mathbf{v}\) 是位置向量且无需注明时，可省略下标，写作 \(\nabla F\)，称为空间梯度、空间散度等。
\end{mydefinition}

\begin{myTheorem}{基础向量微分公式}
令 \(\mathbf{v}\) 为任一向量，\(\mathbf{u}\) 为同方向单位向量，\(\nu\) 为标量大小（\(\mathbf{v} = \nu\mathbf{u}\)），则有：
\begin{align}
\nabla_v \mathbf{v} &= \mathbf{U} \tag{7} \\
\nabla_v \nu &= \mathbf{u} \tag{8} \\
\nabla_v \mathbf{u} &= \nu^{-1}(\mathbf{U} - \mathbf{u}\mathbf{u}) \tag{9}
\end{align}
\end{myTheorem}
\begin{proof}
利用坐标分量对偏导数展开并结合单位并矢 \(\mathbf{U}\) 的定义即可证明上述等式。
\end{proof}

\section{两个粒子的力函数 (Force Function for Two Particles)}

\begin{mydefinition}{引力力函数}
由质量为 \(\overline{m}\) 的粒子 \(\overline{P}\) 作用在质量为 \(m\) 的粒子 \(P\) 上的引力 \(\mathbf{F}\) 可以表示为：
\begin{equation}
\mathbf{F} = \nabla_{\mathbf{p}} V \tag{1}
\end{equation}
其中 \(\mathbf{p}\) 是 \(P\) 相对于 \(\overline{P}\) 的位置矢量，\(V\) 是由下式给出的**力函数**：
\begin{equation}
V = G\overline{m}m p^{-1} + C \tag{2}
\end{equation}
这里 \(p \triangleq (\mathbf{p}^2)^{1/2}\)，\(C\) 为任意常数。
\end{mydefinition}

力函数 \(V\) 总是满足拉普拉斯方程：
\begin{equation}
\nabla^2 V = 0 \tag{4}
\end{equation}
任何满足该方程的解称为**球谐函数**，在此背景下称为**引力势**。

\section{物体和粒子的力函数 (Force Function for a Body and a Particle)}

\begin{myTheorem}{刚体引力力函数}
由质量为 \(\overline{m}\) 的粒子 \(\overline{P}\) 作用在物体 \(B\) 上的合力 \(\mathbf{F}\) 可以写成：
\begin{equation}
\mathbf{F} = \nabla_{\mathbf{R}} V \tag{1}
\end{equation}
其中 \(\mathbf{R}\) 是 \(B^*\) 相对于 \(\overline{P}\) 的位矢，且对于包含 \(N\) 个质点的离散情形：
\begin{equation}
V = G\overline{m}\sum_{i=1}^N m_i p_i^{-1} + C \tag{2}
\end{equation}
对于连续分布：
\begin{equation}
V = G\overline{m}\int p^{-1}\rho d\tau + C \tag{5}
\end{equation}
\(p_i\) 和 \(p\) 分别是 \(\overline{P}\) 到第 \(i\) 个质点和积分微元的距离。
\end{myTheorem}
\begin{proof}
直接对 \(\mathbf{R}\) 求导可发现它正好等于合力 \(\mathbf{F}\)。并且，该势函数也满足拉普拉斯方程 \(\nabla^2 V = 0\)。
\end{proof}

\section{小物体和粒子的力函数 (Force Function for a Small Body and a Particle)}

当距离 \(R\) 远大于物体 \(B\) 的尺寸时，式 (2.11.5) 中的 \(V\) 可以展开为关于 \(1/R\) 的级数：
\begin{equation}
V = \frac{G\overline{m}m}{R}\left[1 + \sum_{i=2}^\infty \nu^{(i)}\right] + C \tag{1}
\end{equation}
其中二阶项为：
\begin{equation}
\nu^{(2)} = \frac{1}{2mR^2}\left[\mathrm{tr}(\mathbf{I}) - 3I_{11}\right] \tag{2}
\end{equation}
这里 \(I_{11}\) 是 \(B\) 关于连接 \(\overline{P}\) 与 \(B^*\) 直线的转动惯量。或者用主转动惯量表示，可得：
\begin{equation}
\nu^{(2)} = \frac{1}{2mR^2}\left[I_1(1 - 3C_{11})^2 + I_2(1 - 3C_{12})^2 + I_3(1 - 3C_{13})^2\right] \tag{6}
\end{equation}
\begin{proof}
将二项式级数展开应用于 \((R^2+2\mathbf{R}\cdot\mathbf{r}+\mathbf{r}^2)^{-1/2}\)，并利用质心与惯性矩的定义进行积分与合并即可。
\end{proof}

\section{用勒让德多项式表示的力函数 (Force Function in Terms of Legendre Polynomials)}

\begin{myTheorem}{勒让德多项式展开}
在式 (2.12.1) 中的项 \(\nu^{(i)}\) 可以通过勒让德多项式 \(P_i\) 表示为：
\begin{equation}
\nu^{(i)} = \frac{1}{m}\int \left(\frac{r}{R}\right)^i P_i(\cos\alpha)\rho d\tau \tag{1}
\end{equation}
其中，\(r\) 是 \(B^*\) 到 \(B\) 上任意点 \(P\) 的距离，\(\alpha\) 是 \(B^*\) 到 \(P\) 与 \(B^*\) 到 \(\overline{P}\) 两直线的夹角。
\end{myTheorem}

勒让德多项式定义为：
\begin{equation}
P_0(y) \triangleq 1 \tag{2}
\end{equation}
且满足生成公式：
\begin{equation}
P_i(y) = \frac{1 \cdot 3 \cdots (2i-1)}{i!}\left[ y^i - \frac{i(i-1)}{(2i-1)\cdot 2}y^{i-2} + \dots \right] \tag{3}
\end{equation}

利用相伴勒让德函数 \(P_i^j\)，并在球坐标 \(r,\gamma,\mu\) 下表达 \(B\) 的质量分布，得到引力势的球谐展开标准形式（国际天文学联合会标准）：
\begin{equation}
V = \frac{G\overline{m}m}{R}\left\{1 + \sum_{i=2}^{\infty}\left[\left(\frac{R_B}{R}\right)^i\sum_{j=0}^i P_i^j(\sin\beta)(C_{ij}\cos j\lambda + S_{ij}\sin j\lambda)\right]\right\} + C \tag{12}
\end{equation}
其中 \(\lambda,\beta,R\) 是 \(\overline{P}\) 的球坐标，\(R_B\) 是归一化参考长度。

对于轴对称物体，取 \(C_{ij}=0, S_{ij}=0\) 且定义 \(J_i \triangleq -C_{i0}\)，则上式化为更简洁的形式：
\begin{equation}
V = \frac{G\overline{m}m}{R}\left[1 - \sum_{i=2}^{\infty}\left(\frac{R_B}{R}\right)^i J_i P_i(\sin\beta)\right] + C \tag{14}
\end{equation}
\begin{proof}
利用生成函数恒等式 \((1 - 2yx + y^2)^{-1/2} = \sum_{i=0}^{\infty} y^i P_i(x)\) 和二项式定理进行展开，并结合球坐标的三角函数正交性质，即可得到上述结果。
\end{proof}

\section{两个小物体的力函数 (Force Function for Two Small Bodies)}

当两个小物体 \(B\) 和 \(\overline{B}\) 间距离 \(R\) 很大时，它们的引力合力可近似为：
\begin{equation}
\widetilde{\mathbf{F}} = -\frac{G\overline{m}m}{R^2}\left[\mathbf{a}_1 + \mathbf{f}^{(2)} + \overline{\mathbf{f}}^{(2)}\right] \tag{1}
\end{equation}
相应的力函数为：
\begin{equation}
\widetilde{V} = \frac{G\overline{m}m}{R}\left[1 + \nu^{(2)} + \overline{\nu}^{(2)}\right] + C \tag{3}
\end{equation}
其中 \(\nu^{(2)}\) 对应 \(B\) 的惯性项，\(\overline{\nu}^{(2)}\) 对应 \(\overline{B}\) 的惯性项。
\begin{proof}
对 \(\widetilde{V}\) 关于 \(\mathbf{R}\) 求导，即可证明 \(\nabla_{\mathbf{R}}\widetilde{V} = \widetilde{\mathbf{F}}\)。
\end{proof}

\section{质心体的力函数 (Force Function for a Centrobaric Body)}

由于质心体的引力效应等效于位于其质心的粒子，因此适用于质心体与任意物体（或粒子）的引力相互作用的力函数，可直接由 2.10 至 2.14 节中的公式通过代入 \(B^*\) 的位矢而获得。

\section{物体与小物体的力函数 (Force Function for a Body and a Small Body)}

\begin{myTheorem}{广义引力力函数}
设 \(V(\mathbf{p})\) 为 \(\overline{B}\) 对位于 \(P\) 处单位质量质点的力函数。当 \(B\) 与 \(\overline{B}\) 的质心距离大于 \(B\) 的最大尺寸时，\(\overline{B}\) 作用在 \(B\) 上的引力合力可近似由力函数 \(\widetilde{V}(\mathbf{R})\) 导出：
\begin{equation}
\widetilde{\mathbf{F}} = \nabla \widetilde{V}(\mathbf{R}) \tag{1}
\end{equation}
其中：
\begin{equation}
\widetilde{V}(\mathbf{R}) = m V(\mathbf{R}) - \frac{1}{2}\mathbf{I}: \nabla\nabla V(\mathbf{R}) + C \tag{2}
\end{equation}
\(m\) 是 \(B\) 的质量，\(\mathbf{I}\) 是 \(B\) 的中心惯性并矢，\(\nabla\nabla V\) 是 \(V(\mathbf{R})\) 的黑塞矩阵。
\end{myTheorem}
\begin{proof}
通过泰勒级数将作用在 \(B\) 上每一点的微元力 \(\nabla_{\mathbf{p}}V\rho d\tau\) 展开，并对 \(B\) 全物体积分，利用质心点定义消去一阶项，利用惯性并矢定义合并二阶项，即可证明上述近似力函数形式。
\end{proof}

\section{粒子施加在物体上的力矩的力函数表达 (Force Function Expressions for the Moment Exerted on a Body by a Particle)}

\begin{myTheorem}{利用力函数计算力矩}
粒子 \(\overline{P}\) 作用在物体 \(B\) 上的引力关于质心 \(B^*\) 的力矩 \(\mathbf{M}\) 可用力函数表示为：
\begin{equation}
\mathbf{M} = -\mathbf{R}\times\nabla V \tag{1}
\end{equation}
如果采用 \(W(\mathbf{a}_1,R) \triangleq V[\mathbf{R}(\mathbf{a}_1,R)]\) 作为 \(\mathbf{a}_1\) 和 \(R\) 的函数，则力矩可写为：
\begin{equation}
\mathbf{M} = -\mathbf{a}_1\times\frac{\partial W}{\partial\mathbf{a}_1} \tag{2}
\end{equation}
对于 \(\widetilde{\mathbf{M}}\)（保留二阶惯性项的近似力矩）：
\begin{equation}
\widetilde{\mathbf{M}} = -\mathbf{a}_1\times\frac{\partial\widetilde{W}}{\partial\mathbf{a}_1} \tag{4}
\end{equation}
其中 \(\widetilde{W} \triangleq -\frac{3G\overline{m}I_{11}}{2R^3} + C\)。
\end{myTheorem}
\begin{proof}
从 \(\mathbf{M} = -\mathbf{R}\times\mathbf{F}\) 出发，利用力函数性质 \(\mathbf{F} = \nabla V\) 和链式法则即可得到上述结果。
\end{proof}

\section{小物体施加在物体上的力矩的力函数表达 (Force Function Expression for the Moment Exerted on a Small Body by a Body)}

当 \(\overline{B}\) 对 \(B\) 的引力力矩需考虑有限大效应时，使用一级近似：
\begin{equation}
\widetilde{\mathbf{M}} = -\mathbf{I} \times \nabla\nabla V(\mathbf{R}) \tag{1}
\end{equation}
其中 \(\times\) 表示并矢的**叉积**运算。
\begin{proof}
将 \(B\) 中质点的引力场按泰勒级数展开，由于质心项积分为零，通过保留剩余一阶项并结合惯性并矢 \(\mathbf{I}\) 的定义即可得到。注意 \(\nabla\nabla V\) 的对称性会使运算得到化简。
\end{proof}